\documentclass[12pt, a4paper]{article}

% ── Packages ──────────────────────────────────────────────────────────────────
\usepackage[a4paper, margin=2.5cm]{geometry}
\usepackage{graphicx}
\usepackage{float}
\usepackage{amsmath}
\usepackage{amssymb}
\usepackage{xcolor}
\usepackage{mdframed}
\usepackage{parskip}
\usepackage{enumitem}
\usepackage{titlesec}
\usepackage{booktabs}
\usepackage{hyperref}
\usepackage[T1]{fontenc}
\usepackage{tikz}
\usetikzlibrary{arrows.meta}

% ── Graphics path ─────────────────────────────────────────────────────────────
\graphicspath{{../../Images_for_latex/Lecture_4/}}

% ── Colours ───────────────────────────────────────────────────────────────────
\definecolor{keyblue}{RGB}{30, 80, 160}
\definecolor{keybluebg}{RGB}{220, 235, 255}
\definecolor{noteorange}{RGB}{200, 100, 0}
\definecolor{noteorangebg}{RGB}{255, 240, 215}
\definecolor{warnred}{RGB}{180, 30, 30}
\definecolor{warnredbg}{RGB}{255, 220, 220}
\definecolor{defgreen}{RGB}{20, 110, 50}
\definecolor{defgreenbg}{RGB}{215, 245, 225}
\definecolor{headernavy}{RGB}{25, 35, 90}

% ── Box environments ──────────────────────────────────────────────────────────
\newmdenv[linecolor=keyblue,backgroundcolor=keybluebg,linewidth=1.5pt,
  leftline=true,rightline=false,topline=false,bottomline=false,
  innerleftmargin=10pt,innerrightmargin=8pt,
  innertopmargin=6pt,innerbottommargin=6pt,
  skipabove=8pt,skipbelow=8pt]{keybox}

\newmdenv[linecolor=noteorange,backgroundcolor=noteorangebg,linewidth=1.5pt,
  leftline=true,rightline=false,topline=false,bottomline=false,
  innerleftmargin=10pt,innerrightmargin=8pt,
  innertopmargin=6pt,innerbottommargin=6pt,
  skipabove=8pt,skipbelow=8pt]{notebox}

\newmdenv[linecolor=warnred,backgroundcolor=warnredbg,linewidth=1.5pt,
  leftline=true,rightline=false,topline=false,bottomline=false,
  innerleftmargin=10pt,innerrightmargin=8pt,
  innertopmargin=6pt,innerbottommargin=6pt,
  skipabove=8pt,skipbelow=8pt]{warnbox}

\newmdenv[linecolor=defgreen,backgroundcolor=defgreenbg,linewidth=1.5pt,
  leftline=true,rightline=false,topline=false,bottomline=false,
  innerleftmargin=10pt,innerrightmargin=8pt,
  innertopmargin=6pt,innerbottommargin=6pt,
  skipabove=8pt,skipbelow=8pt]{defbox}

% ── Section formatting ────────────────────────────────────────────────────────
\titleformat{\section}[block]{\large\bfseries\color{headernavy}}{}{0pt}{}
  [\vspace{2pt}\rule{\linewidth}{0.8pt}\vspace{4pt}]
\titleformat{\subsection}[block]{\normalsize\bfseries\color{keyblue}}{}{0pt}{}

% ── Document ──────────────────────────────────────────────────────────────────
\begin{document}

\begin{center}
  {\LARGE\bfseries\color{headernavy} Sensors \& Systems}\\[4pt]
  {\large Kalman Filter for Nonlinear Systems --- EKF \& UKF}\\[2pt]
  {\normalsize Torben Knudsen \quad|\quad Electronic Systems, 8th Semester}
\end{center}

\texttt{np.linalg.eigh($(n+\lambda)P$)} decomposes the scaled covariance matrix into:
\begin{itemize}[leftmargin=1.4em]
    \item \texttt{eigvals} $= \Lambda$ --- how much variance exists in each direction
    \item \texttt{eigvecs} $= V$ --- which directions the variance points in
\end{itemize}
Think of it as finding the natural axes of the uncertainty ellipse and how stretched it is along each axis.

\texttt{S = eigvecs @ np.diag(np.sqrt(eigvals))} builds the matrix square root
\[
  S = V\sqrt{\Lambda}
\]
so each column of $S$ points in an eigenvector direction (shape of uncertainty)
scaled by $\sqrt{\lambda_i}$ (how far to push in that direction).
Together $S \approx \sqrt{(n+\lambda)P}$, giving the sigma point offsets from the mean.

\begin{notebox}
  \texttt{np.maximum(eigvals, 0)} clips any negative eigenvalues to zero before
  taking the square root. $P$ should be positive definite by construction, so all
  eigenvalues should be $> 0$, but floating-point accumulation over many steps can
  produce tiny negative values. Taking $\sqrt{\text{negative}}$ would give \texttt{NaN}
  and silently break the filter --- the clip prevents this.
\end{notebox}
\end{document}
